\section{Setting and results}
\label{sec:theory}

\subsection{Two distances}

Let $(X,d)$ be a metric space of inputs and $f:X\to\R$ a function. A data point
is a pair $(x,y)\in X\times\R$. We give $X\times\R$ the Euclidean product metric
$\rho\bigl((x,y),(x',y')\bigr)=\sqrt{d(x,x')^2+(y-y')^2}$ and write
$\Gamma_f=\{(x,f(x)):x\in X\}$ for the graph of $f$, the ``line'' of the
conjecture. Two quantities compare a data point with the line:
\begin{align*}
  \vd_f(x,y) &= |y-f(x)| && \text{(vertical error)},\\
  \nd_f(x,y) &= \inf_{x'\in X}\rho\bigl((x,y),(x',f(x'))\bigr) && \text{(nearest-point distance)}.
\end{align*}
Regression losses are functions of the residual $y-f(x)$, so they see $\vd_f$.
The conjecture is stated with $\nd_f$.

A \emph{random function} $F$ is a probability space $(\Theta,\pi)$ of parameters
together with functions $f_\theta:X\to\R$ such that $\theta\mapsto f_\theta(x)$
is measurable for every $x$. Its \emph{one-point marginal} at $x$ is the law
$\nu^F_x$ of $f_\theta(x)$ under $\pi$. For a cost $c$ and a data distribution
$\mu$ on $X\times\R$, the expected cost at initialisation is
$\E_{(x,y)\sim\mu}\,\E_{\theta\sim\pi}\,c(f_\theta;x,y)$.\footnote{We assume the
inner expectation is a measurable function of $(x,y)$. This holds if
$(\theta,x)\mapsto f_\theta(x)$ is jointly measurable, and always when $\mu$ is
the empirical measure of a finite data set.} For a family of random functions
indexed by a size $m$, ``a bigger model starts closer'' means that this
expectation is strictly decreasing in $m$. Claim (G) uses the cost $\nd$ and
claim (L) uses costs of the form $\ell(y-f(x))$. We make ``squiggly'' concrete
by the frequency of a wave of fixed amplitude and, for networks, by arc length
and the number of turns.

\subsection{The slope sandwich}

\begin{theorem}[Slope sandwich]\label{thm:sandwich}
If $f$ is $K$-Lipschitz, then for every data point $(x,y)$
\[
  \frac{\vd_f(x,y)}{\sqrt{1+K^2}}\;\le\;\nd_f(x,y)\;\le\;\vd_f(x,y).
\]
\end{theorem}

\begin{proof}
The point $(x,f(x))$ lies on $\Gamma_f$ at distance $|y-f(x)|$ from $(x,y)$,
which gives the upper bound. For the lower bound take any $x'\in X$ and put
$a=d(x,x')$ and $b=|y-f(x')|$. By the triangle inequality and the Lipschitz
bound, $|y-f(x)|\le|y-f(x')|+|f(x')-f(x)|\le b+Ka$. By the Cauchy--Schwarz
inequality, $b+Ka\le\sqrt{1+K^2}\,\sqrt{a^2+b^2}=\sqrt{1+K^2}\;\rho\bigl((x,y),(x',f(x'))\bigr)$.
Taking the infimum over $x'$ gives $\vd_f(x,y)\le\sqrt{1+K^2}\,\nd_f(x,y)$.
\end{proof}

Both bounds are attained: for $X=\R$ and $f(x)=Kx$, the point $(0,y)$ has
$\nd_f=|y|/\sqrt{1+K^2}$, and for a constant $f$ the two distances coincide.
Taking expectations, the same two inequalities hold between the mean distances
of any random $K$-Lipschitz function and any data distribution.

The theorem says what a small nearest-point distance means. If
$\nd_f(x,y)<\eta$, there is an input $x'$ with $d(x,x')<\eta$ and
$|y-f(x')|<\eta$: a perturbation of the \emph{input} smaller than $\eta$ brings
the output within $\eta$ of the label. When $\nd_f$ is small and $\vd_f$ is not,
the function is steep near $x$: its Lipschitz constant is at least
$\sqrt{(\vd_f/\nd_f)^2-1}$. That describes a function that is
sensitive to its input near the data, not one that fits the data.

\subsection{The marginal lemma}

\begin{theorem}[Marginal lemma]\label{thm:marginal}
Let $F$ be a random function and $\ell:\R\to[0,\infty)$ a measurable loss. For
every data point $(x,y)$,
\[
  \E_{\theta}\,\ell\bigl(y-f_\theta(x)\bigr)=\int_{\R}\ell(y-t)\,\nu^F_x(dt).
\]
Consequently, if two random functions $F$ and $G$ have $\nu^F_x=\nu^G_x$ for
every $x$, their expected losses are equal for every such $\ell$ and every data
distribution $\mu$.
\end{theorem}

\begin{proof}
By definition $\nu^F_x$ is the image of $\pi$ under $\theta\mapsto f_\theta(x)$.
The displayed identity is the change-of-variables formula for an image measure,
applied to the measurable function $t\mapsto\ell(y-t)$. Integrating both sides
against $\mu$ gives the second statement.
\end{proof}

The proof is one line, and the content is in what is absent: the joint law of
$\bigl(f_\theta(x),f_\theta(x')\bigr)$ at two inputs does not enter. How $f$
moves between inputs is what ``squiggly'' means, so the expected initial loss
cannot register squiggliness. The lemma is about the expectation over $\theta$
only. The variance of the initial loss, the probability of a low initial loss,
and everything that happens during training do depend on joint laws.

\begin{corollary}[Random-phase waves]\label[corollary]{cor:phase}
Let $g:\R\to\R$ be measurable and $T$-periodic, let $\varphi$ be uniform on
$[0,T)$, and for a frequency $w\in\R$ let $f_{w,\varphi}(x)=g(wx+\varphi)$ on
$X=\R$. Then the one-point marginal of $f_{w,\cdot}$ at every $x$ is the law of
$g(\varphi)$, for every $w$. Hence the expected loss
$\E_{(x,y)\sim\mu}\E_\varphi\,\ell\bigl(y-f_{w,\varphi}(x)\bigr)$ is the same for
every $w$, every measurable $\ell\ge0$ and every $\mu$. No sequence of
frequencies makes it decrease.
\end{corollary}

\begin{proof}
The uniform distribution on the circle $\R/T\mathbb{Z}$ is invariant under
translation, so $wx+\varphi \bmod T$ is uniform on $[0,T)$ for every fixed $w$
and $x$, and $g(wx+\varphi)$ has the law of $g(\varphi)$. Apply
\cref{thm:marginal}.
\end{proof}

For the curves of \cref{sec:experiments}, $g=p\sin$ with peak $p\le1$, and data
with $y$ uniform on $[-1,1]$ independent of $x$. Since
$\E_y|y-s|=(1+s^2)/2$ for $|s|\le1$, the expected vertical error is
$(1+p^2/2)/2$ at every frequency.

\begin{corollary}[ReLU networks under He initialisation]\label[corollary]{cor:he}
Let $f_\theta(x)=u^\top\relu\bigl(W^L\relu(\cdots\relu(W^1x))\bigr)$ for
$x\in\R^{d}$, with $\relu$ the ReLU, no biases, hidden widths $n_1,\dots,n_L$,
$n_0=d$, and independent weights $W^k_{ij}\sim N(0,2/n_{k-1})$ and
$u_i\sim N(0,2/n_L)$. Then for every $x$ and $y$,
\[
  \E f_\theta(x)=0,\qquad \E f_\theta(x)^2=\frac{2\|x\|^2}{d},\qquad
  \E\bigl(y-f_\theta(x)\bigr)^2=y^2+\frac{2\|x\|^2}{d},
\]
whatever the depth and the widths. If all weights are instead $N(0,1)$, the last
quantity is $y^2+\|x\|^2\prod_{k=1}^{L}(n_k/2)$.
\end{corollary}

The proof is a layer-by-layer recursion for $\E\|h^k\|^2$ and is in
\cref{app:proofs}. It uses only that the weights are independent, symmetric and
of the stated variance; preserving the second moment across layers is the
design criterion of this initialisation \citep{he2015delving}. In the terms of
\cref{thm:marginal}, the squared loss reads only the first two moments of
$\nu_x$, and these do not depend on the width. The full law $\nu_x$ does depend
on the width, so for other losses the corollary gives no exact statement.

\subsection{The tube identity and a bound from length}

For a non-empty set $S$ in a metric space, write
$N_r(S)=\{q:\dist(q,S)\le r\}$ for the closed tube of radius $r$ around $S$.

\begin{theorem}[Tube identity]\label{thm:tube}
Let $S$ be a non-empty subset of a metric space and $P$ a random point of that
space. Then
\[
  \E\,\dist(P,S)=\int_0^\infty\P\bigl(\dist(P,S)>r\bigr)\,dr
  =\int_0^\infty\P\bigl(P\notin N_r(S)\bigr)\,dr ,
\]
where both sides may be infinite.
\end{theorem}

\begin{proof}
The function $q\mapsto\dist(q,S)$ is $1$-Lipschitz, hence measurable, and
non-negative. For a non-negative random variable $Z$ we have
$Z=\int_0^\infty\mathbf 1\{Z>r\}\,dr$, and Tonelli's theorem gives
$\E Z=\int_0^\infty\P(Z>r)\,dr$. This is the layer-cake formula
\citep[Section~1.13]{lieb2001analysis}. Take $Z=\dist(P,S)$.
\end{proof}

The identity holds for any set in any metric space. It reduces the mean
distance to the volume of tubes, which for curves in the plane is controlled by
length.

\begin{lemma}\label[lemma]{lem:tube-area}
Let $S\subset\R^2$ be the image of a rectifiable curve of length $l$. Then the
area of $N_r(S)$ is at most $2lr+\pi r^2$ for every $r\ge0$.
\end{lemma}

The proof, by polygonal approximation, is in \cref{app:proofs}.

\begin{proposition}[Length bound]\label[proposition]{prop:length}
Let $\Omega\subset\R^2$ be measurable with area $V\in(0,\infty)$, let $P$ be
uniform on $\Omega$, and let $S$ be the image of a rectifiable curve of length
$l$. Then
\[
  \E\,\dist(P,S)\;\ge\;\frac{V}{4l+2\sqrt{\pi V}} .
\]
\end{proposition}

\begin{proof}
Put $r_0=V/(2l+\sqrt{\pi V})$. For $0\le r\le r_0$ we have
$\pi r\le\pi V/\sqrt{\pi V}=\sqrt{\pi V}$, so by \cref{lem:tube-area} the area
of $N_r(S)$ is at most $(2l+\sqrt{\pi V})\,r=Vr/r_0$. Therefore
$\P(P\notin N_r(S))\ge1-r/r_0$ for $r\le r_0$, and by \cref{thm:tube}
$\E\,\dist(P,S)\ge\int_0^{r_0}(1-r/r_0)\,dr=r_0/2$.
\end{proof}

To leading order in $l$ the bound is $V/(4l)$. We draw three consequences.
Length is necessary: a mean distance of at most $\varepsilon$ requires
$l\ge V/(4\varepsilon)-\sqrt{\pi V}/2$. Length is not sufficient: a part of the
curve further than $r$ from $\Omega$ adds nothing to $N_r(S)\cap\Omega$, so
length outside the data region is wasted (\cref{sec:line}). And the bound is
nearly attained only if the tube around the curve stays inside $\Omega$ and
overlaps itself little, that is, if the length is spread evenly over the
region. The proposition is for curves in the plane; for inputs of dimension
above one the graph is a surface, and its area takes the role of length.

\subsection{The sine family}

The experiments use $f_{w,\varphi}(x)=p\sin(w\pi x+\varphi)$ on $x\in[-1,1]$
with $\varphi$ uniform on $[0,2\pi)$, and a point $P$ uniform on
$\Omega=[-1,1]^2$, so $V=4$. The curve makes $w$ full periods across the box and
is $p\,w\pi$-Lipschitz. Write $D_p(w)$ for the mean of $\nd$ over $P$ and
$\varphi$, and $l(w)$ for the arc length.

\begin{proposition}[The sine that spans the data; leading order]\label[proposition]{prop:sine}
For $p=1$, as $w\to\infty$,
\[
  \frac{l(w)}{w}\to4,\qquad
  w\,D_1(w)\to\frac{2-8/\pi^2}{4}\approx0.2974,\qquad
  D_1(w)\,l(w)\to2-\frac{8}{\pi^2}\approx1.189 .
\]
\end{proposition}

The first limit is proved in \cref{app:proofs}. The second is a leading-order
derivation. At height $y$ the curve crosses on two interleaved lattices of
period $2/w$, with gaps $(1/w)(1\mp2\arcsin(y)/\pi)$. At high frequency the
strokes are nearly vertical, so the nearest point of the curve is, to leading
order, the nearest crossing at the same height. Averaging that distance over
$x$ and then over $y$, with $\int_0^1\arcsin(y)^2\,dy=\pi^2/4-2$, gives the
constant. The step from horizontal to Euclidean distance is argued in outline
in \cref{app:proofs}, without all details, and the constant is checked
numerically in \cref{sec:saturation}. The limit $1.189$ exceeds the value
$V/4=1$ of \cref{prop:length} because the strokes of a sine are not evenly
spaced.

\begin{proposition}[Band-edge floor]\label[proposition]{prop:floor}
For $0<p\le1$ and $w\ge1$,
\[
  \frac{(1-p)^2}{2}\;\le\;D_p(w)\;\le\;\frac{(1-p)^2}{2}+\frac{2}{w}.
\]
\end{proposition}

The proof is in \cref{app:proofs}. The curve stays in the band $|y|\le p$, so a
point at height $|y|>p$ is never nearer than $|y|-p$, and
$\E\max(0,|y|-p)=(1-p)^2/2$. Raising the frequency removes the rest at rate
$O(1/w)$. The floor depends on the amplitude relative to the range of the data.
It does not depend on the number of data points: $D_p(w)$ is the mean of a
per-point quantity, so its expectation is the same for a sample of any size.

\subsection{The covering bound and the theorem for sines}
\label{sec:cover}

The upper bound in \cref{prop:floor} is an instance of a general fact, which
says what a function must do to be near the data.

\begin{definition}[Covering]\label[definition]{def:cover}
A function $f:X\to\R$ \emph{covers} a set of labels $S\subset\R$ \emph{within}
$\rho$ if for every $x\in X$ and every $y\in S$ there is an $x'\in X$ with
$d(x,x')\le\rho$ and $f(x')=y$.
\end{definition}

The smallest such $\rho$ is an oscillation scale: how far one must move in the
input to see every label.

\begin{theorem}[Covering bound]\label{thm:cover}
If $f$ covers $S$ within $\rho$, then $\nd_f(x,y)\le\rho$ for every data point
with $y\in S$. If in addition $f$ is $K$-Lipschitz, then
$\vd_f(x,y)\le\sqrt{1+K^2}\,\rho$ for every such point.
\end{theorem}

\begin{proof}
Take $x'$ as in \cref{def:cover}. The point $(x',y)=(x',f(x'))$ lies on
$\Gamma_f$ at distance $d(x,x')\le\rho$ from $(x,y)$. The second statement
follows from \cref{thm:sandwich}.
\end{proof}

\begin{corollary}[Nearness is a distance to move]\label[corollary]{cor:warp}
If $f$ covers $S$ within $\rho$, then for every target $g:X\to S$ there is a map
$h:X\to X$ with $d(x,h(x))\le\rho$ for every $x$ and $f\circ h=g$.
\end{corollary}

\begin{proof}
For each $x$ choose $h(x)$ as the point $x'$ of \cref{def:cover} for the label
$g(x)$.
\end{proof}

The corollary says in what sense the nearness can be used. The loss at
initialisation does not see it (\cref{thm:marginal}). But a function that
covers the labels within $\rho$ fits any data with labels in $S$ once each
input is moved by at most $\rho$. For a single data point the movement is a
shift, which for a wave is a change of phase. For several points the shifts
differ, and $h$ is a warp of the input, which in a network is what the earlier
layers can supply. The corollary does not say that $h$ is continuous, or that
training finds it.

So a small oscillation scale is sufficient for a small nearest-point distance,
and by the second statement of \cref{thm:cover} it has a price: a function that covers labels at
vertical distance $\vd$ within $\rho$ has slope at least
$\sqrt{(\vd/\rho)^2-1}$. Claims (G) and (P) for a family of models therefore
come down to how its oscillation scale and its slope depend on its size. For
sines both are explicit.

The sine is one case of a statement about any function. For a function $f$ on
a normed space $X$ and $w>0$, write $f_w(x)=f(wx)$ for $f$ squeezed sideways by
the factor $w$.

\begin{theorem}[The squiggly line theorem]\label{thm:general}
Let $F$ be a random function on a normed space $X$ such that every draw
$f_\theta$ covers $S$ within $\rho$ and is $K$-Lipschitz, and such that the
one-point marginal $\nu^F_x$ is the same at every $x$. Let $F_w$ be $F$ with
every draw squeezed by $w>0$.
\begin{enumerate}\setlength{\itemsep}{0pt}
\item[(i)] \emph{Closer and steeper.} For every draw and every data point with
$y\in S$,
$\vd/\sqrt{1+(wK)^2}\le\nd\le\rho/w$.
\item[(ii)] For a probability measure $\mu$ on data points with $y\in S$ almost
surely, the expected nearest-point distance of $F_w$ is at most $\rho/w$ and
tends to $0$ as $w\to\infty$.
\item[(iii)] \emph{No better at the start.} For every measurable loss
$\ell\ge0$, the expected loss of $F_w$ is the same for every $w$.
\item[(iv)] For every target $g:X\to S$ and every draw there is a map $h$ with
$d(x,h(x))\le\rho/w$ and $f_{\theta,w}\circ h=g$.
\end{enumerate}
\end{theorem}

\begin{proof}
If $f$ covers $S$ within $\rho$ then $f_w$ covers $S$ within $\rho/w$: given
$x$ and $y$, take $z$ with $\|wx-z\|\le\rho$ and $f(z)=y$, and put $x'=z/w$.
If $f$ is $K$-Lipschitz then $f_w$ is $wK$-Lipschitz. So (i) is
\cref{thm:sandwich,thm:cover}, (ii) follows from (i), and (iv) is
\cref{cor:warp}. For (iii), the marginal of $F_w$ at $x$ is $\nu^F_{wx}$, which
does not depend on $w$ or $x$; apply \cref{thm:marginal}.
\end{proof}

No part of the statement uses the dimension of $X$ or a period. For the sine
the constants are explicit.

\begin{theorem}[The squiggly line theorem, for sines]\label{thm:sine}
Let $X=\R$, $p>0$ and $w>0$, and let $f_{w,\varphi}(x)=p\sin(w\pi x+\varphi)$.
\begin{enumerate}\setlength{\itemsep}{0pt}
\item[(i)] For every phase $\varphi$ and every data point $(x,y)$ with
$|y|\le p$,
\[
  \frac{\vd_{f_{w,\varphi}}(x,y)}{\sqrt{1+(p\,w\pi)^2}}\;\le\;\nd_{f_{w,\varphi}}(x,y)\;\le\;\frac1w .
\]
\item[(ii)] If $\mu$ is a probability measure on data points with $|y|\le p$
almost surely and $\varphi$ is uniform on $[0,2\pi)$, the expected
nearest-point distance is at most $1/w$, and tends to $0$ as $w\to\infty$.
\item[(iii)] For every measurable loss $\ell\ge0$ and every $\mu$, the expected
loss $\E_{(x,y)\sim\mu}\E_\varphi\,\ell\bigl(y-f_{w,\varphi}(x)\bigr)$ is the
same for every $w$.
\end{enumerate}
\end{theorem}

\begin{proof}
(i) The function is $p\,w\pi$-Lipschitz, which gives the lower bound by
\cref{thm:sandwich}. For the upper bound we show that it covers $[-p,p]$ within
$1/w$. Given $x$ and $y$ with $|y|\le p$, put $t=\arcsin(y/p)$ and
$\theta_0=w\pi x+\varphi$, let $k$ be an integer nearest to
$(\theta_0-t)/(2\pi)$, and put $\theta=t+2\pi k$. Then $|\theta-\theta_0|\le\pi$.
The input $x'=(\theta-\varphi)/(w\pi)$ satisfies $|x-x'|=|\theta-\theta_0|/(w\pi)\le1/w$
and $f_{w,\varphi}(x')=p\sin t=y$. Apply \cref{thm:cover}.
(ii) The bound in (i) holds for every phase and almost every data point, and
$\nd\ge0$. (iii) is \cref{cor:phase}.
\end{proof}

Parts (i) and (ii) are claim (G), with a rate: both ends of the sandwich fall
like $1/w$. Part (iii) is the failure of claim (L). By \cref{cor:warp}, any
target with values in $[-p,p]$ is fitted exactly by $f_{w,\varphi}$ after a warp
of the input of size at most $1/w$. For one data point this is a change of
phase of at most $\pi$, at every frequency: a squigglier sine is closer in the
input and exactly as far in the phase. On the bounded domain
$[-1,1]$ of the experiments the same argument gives $2/w$ in place of $1/w$,
which is the upper bound of \cref{prop:floor}; the exact constant is the one in
\cref{prop:sine}.
